\section{Рабочий проект}

Можно выделить несколько модулей, использованных при разработке
виртуальной файловой системы: командный интепретатор, интерфейс виртуальной файловой системы, главная загрузочная запись, файловая система FAT-32, файловая система ISO-9660(CDFS), виртуальная файловая система, работа с блоками диска, драйвер ATA.

\subsection{Модуль "<shell.lsp">}

Функции модуля интерфейса виртуальной файловой системы представлены в таблице \ref{cmd:table}.

\renewcommand{\arraystretch}{0.8} % уменьшение расстояний до сетки таблицы
\begin{xltabular}{\textwidth}{|p{3.5cm}|p{3.5cm}|>{\setlength{\baselineskip}{0.7\baselineskip}}X|}
	\caption{Описание функций, используемых в модуле командного интерпретатора\label{cmd:table}}\\
	\hline \centrow \setlength{\baselineskip}{0.7\baselineskip} Название функции & \centrow \setlength{\baselineskip}{0.7\baselineskip} Аргументы & \centrow Описание функции \\
	\hline \centrow 1 & \centrow 2 & \centrow 3\\ \hline
	\endfirsthead
	\caption*{Продолжение таблицы \ref{cmd:table}}\\
	\hline \centrow 1 & \centrow 2 & \centrow 3\\ \hline
	\finishhead
	key-down-handler & scan & Обработчик нажатия клавиш с клавиатуры, где \texttt{scan} -- скан код нажатой клавиши. \\ \hline
	process-command & com & Обработка введённой команды, где \texttt{com} -- список введённых символов. \\ \hline
	print-start & Аргументов нет & Вывод на экран строки активности командного интерпретатора. \\ \hline
	parse-commands & Аргументов нет & Парсер всех возможных команд командного интерпретатора. Выполняет соответствующую команду. \\ \hline
	parse-command & com fun & Парсер определённый команды, где \texttt{com} -- строка названия команды, \texttt{fun} -- функция, которую выполняет команда. \\ \hline
	parse-elem-str & str & Предикат проверки корректности названия введённой команды, где \texttt{str} -- строка названия команды.  \\ \hline
	get-str & Аргументов нет & Получить строку названия команды. \\ \hline
	is-com-separator & c & Предикат проверки символа c на принадлежность к разделителям. \\ \hline
	cat-command-action & args & Выполнение команды вывода содержимого файла на экран с аргументами \texttt{args}(путь). \\ \hline
	copy-command-action & args & Выполнение команды копирования файла с аргументами \texttt{args}(путь-откуда путь-куда). \\ \hline
	append-command-action & args & Выполнение команды добавление текста в конец файла с аргументами \texttt{args}(путь текст). \\ \hline
	chattr-command-action & args & Выполнение команды изменения атрибутов файла или каталога с аргументами \texttt{args}(путь атрибуты).
\end{xltabular}
\renewcommand{\arraystretch}{1.0} % восстановление сетки

\subsection{Модуль "<vfs.lsp">}

Функции модуля интерфейса виртуальной файловой системы представлены в таблице \ref{vfs:table}.

\renewcommand{\arraystretch}{0.8} % уменьшение расстояний до сетки таблицы
\begin{xltabular}{\textwidth}{|p{3.5cm}|p{3.5cm}|>{\setlength{\baselineskip}{0.7\baselineskip}}X|}
	\caption{Описание функций, используемых в модуле интерфейса виртуальной файловой системы\label{vfs:table}}\\
	\hline \centrow \setlength{\baselineskip}{0.7\baselineskip} Название функции & \centrow \setlength{\baselineskip}{0.7\baselineskip} Аргументы & \centrow Описание функции \\
	\hline \centrow 1 & \centrow 2 & \centrow 3\\ \hline
	\endfirsthead
	\caption*{Продолжение таблицы \ref{vfs:table}}\\
	\hline \centrow 1 & \centrow 2 & \centrow 3\\ \hline
	\finishhead
	cur-dir & Аргументов нет & Получение пути текущего каталога. \\ \hline
	list-dir & path & Получение содержимого каталога по пути \texttt{path}. \\ \hline
	create-dir & path & Создание каталога по пути \texttt{path} где последний элемент пути это название нового каталога. \\ \hline
	remove-dir & path & Удаление каталога по пути \texttt{path} если он пуст. \\ \hline
	change-dir & path & Перемещение в каталог по пути \texttt{path}. \\ \hline
	create-file & path & Создание файла по пути \texttt{path} где последний элемент пути это название нового файла.  \\ \hline
	remove-file & path & Удаление файла по пути \texttt{path}. \\ \hline
	open-file & path & Открытие файла как поток для чтения и записи по пути \texttt{path} и вернуть файловый объект. \\ \hline
	is-directory & path & Проверка файла по пути \texttt{path} на каталог. \\ \hline
	rename & path name & Переименовывание файла или каталога по пути \texttt{path} на name. \\ \hline
	free-space & Аргументов нет & Получение объёма свободного места в разделе. \\ \hline
	fstat & path & Получение метаданных файла или каталога. \\ \hline
	set-attr & path\newline new-attributes & Изменение атрибутов файла или каталога по пути \texttt{path} на \texttt{new-attributes}. \\ \hline
	find-entry & dir name & Поиск файла или каталога с названием name в каталоге \texttt{dir}. \\ \hline
	parse-path & str-path & Преобразование строки пути \texttt{str-path} в список пути. \\ \hline
	parse-up-dir & path-list & Преобразовать список пути \texttt{path-list} убрав все \textquotedbl..\textquotedbl. \\ \hline
	load-path & path & Проверяет и при необходимости загружает путь \texttt{path}. Возвращает файловый объект файла или каталога или вызывает ошибку если пути нет. \\ \hline
	get-parent-dir & path & Получение файлового объекта родительского каталога файла или каталога по строке пути \texttt{path}.
\end{xltabular}
\renewcommand{\arraystretch}{1.0} % восстановление сетки

\subsection{Модуль "<mbr.lsp">}

Функции модуля главной загрузочной записи представлены в таблице \ref{mbr:table}.

\renewcommand{\arraystretch}{0.8} % уменьшение расстояний до сетки таблицы
\begin{xltabular}{\textwidth}{|p{3.5cm}|p{3.5cm}|>{\setlength{\baselineskip}{0.7\baselineskip}}X|}
	\caption{Описание функций, используемых в модуле главной загрузочной записи\label{mbr:table}}\\
	\hline \centrow \setlength{\baselineskip}{0.7\baselineskip} Название функции & \centrow \setlength{\baselineskip}{0.7\baselineskip} Аргументы & \centrow Описание функции \\
	\hline \centrow 1 & \centrow 2 & \centrow 3\\ \hline
	\endfirsthead
	\caption*{Продолжение таблицы \ref{mbr:table}}\\
	\hline \centrow 1 & \centrow 2 & \centrow 3\\ \hline
	\finishhead
	parse-partition-rec & Аргументов нет & Парсер записи таблицы раздеров в MBR. \\ \hline
	load-partition & disk-num part-num & Загрузить файловую систему из раздела \texttt{part-num} с диска \texttt{disk-num}.
\end{xltabular}
\renewcommand{\arraystretch}{1.0} % восстановление сетки

\subsection{Модули "<fat32-class.lsp">, "<fat32-fat.lsp">, "<fat32-entry.lsp">, "<fat32-fs.lsp">, "<fat32-file.lsp">}

Функции модуля файловой системы FAT-32 представлены в таблице \ref{fat32:table}.

\renewcommand{\arraystretch}{0.8} % уменьшение расстояний до сетки таблицы
\begin{xltabular}{\textwidth}{|p{3.5cm}|p{3.5cm}|>{\setlength{\baselineskip}{0.7\baselineskip}}X|}
	\caption{Описание функций, используемых в модуле файловой системы FAT-32\label{fat32:table}}\\
	\hline \centrow \setlength{\baselineskip}{0.7\baselineskip} Название функции & \centrow \setlength{\baselineskip}{0.7\baselineskip} Аргументы & \centrow Описание функции \\
	\hline \centrow 1 & \centrow 2 & \centrow 3\\ \hline
	\endfirsthead
	\caption*{Продолжение таблицы \ref{fat32:table}}\\
	\hline \centrow 1 & \centrow 2 & \centrow 3\\ \hline
	\finishhead
	parse-bios-parameter-block & Аргументов нет & Парсер блока параметров BIOS. \\ \hline
	parse-fsinfo & Аргументов нет & Парсер структуры FSInfo. \\ \hline
	update-fsinfo & free-count free-num & Обновить значения структуры fsinfo на диске на \texttt{free-count} и \texttt{free-num}. \\ \hline
	fat-elem-pos & idx & Получить пару из смещение сектора + смещение в секторе для элемента с индексом \texttt{idx} таблицы FAT. \\ \hline
	fat-read-sector & fat-table sector-num & Прочитать сектор \texttt{sector-num} таблицы FAT, записав содержимое сектора таблицы в массив \texttt{fat-table}. \\ \hline
	fat-check-for-read & fat-table fat-read-sectors elem-num & Проверить прочитан ли сектор таблицы FAT, в котором находится элемент \texttt{elem-num} через массив прочитанных секторов \texttt{fat-read-sectors}. Если нет то прочитать сектор таблицы FAT в \texttt{fat-table}. \\ \hline
	get-fat-chain & block-num & Получить цепочку блоков из таблицы FAT, начиная с блока \texttt{block-num}. \\ \hline
	get-free-block & fat-table fat-read-sectors free-block-num & Найти первый свободный блок в таблице FAT \texttt{fat-table}, где массив прочитанных секторов таблицы FAT \texttt{fat-read-sectors}, начиная поиск с \texttt{free-block-num}. \\ \hline
	update-fat-elem & idx val & Установить элемент с индексом \texttt{idx} в таблице FAT на диске на значение \texttt{val} при необходимости обновить во всех копиях таблицы FAT. \\ \hline
	append-fat-chain & first-block new-idx & Добавить элемент \texttt{new-idx} в цепочку блоков с первым блоком \texttt{first-block}. \\ \hline
	fat32-is-special-name & name & Проверка на специальное имя каталога. \\ \hline
	fat32-get-date-time & date-bytes time-bytes creation-time-ms & Получить список даты и времени из массивов байт \texttt{date-bytes}, \texttt{time-bytes} и \texttt{creation-time-ms}. \\ \hline
	fat32-generate-date-time & date-time & Получить массив байт для даты и времени для записи в каталоге из списка даты и времени \texttt{date-time}. \\ \hline
	fat32-get-attributes & attribute-byte & Получить список атрибутов из записи в каталоге из байта \texttt{attribute-byte}. \\ \hline
	fat32-generate-attribute-byte & attributes & Получить байт атрибутов для записи в каталоге из списка атрибутов \texttt{attributes}. \\ \hline
	parse-lfn-entry & Аргументов нет & Парсер структуры записи длинного имени. \\ \hline
	parse-fat32-dir-entry & Аргументов нет & Парсер структур записи в каталоге. \\ \hline
	fat32-calc-lfn-count & name & Посчитать количество записей длинного имени для имени \texttt{name}. \\ \hline
	fat32-generate-checksum & short-name & Посчитать контрольную сумму короткого имени для записи длинного имени. \\ \hline
	fat32-create-lfn-entry & name-part-str lfn-num checksum & Создать одну запись длинного имени для строки \texttt{name-part-str}, с номером \texttt{lfn-num} и контрольной суммой \texttt{checksum}. \\ \hline
	fat32-create-lfn-entries & entry & Создать записи длинного имени для записи в каталоге \texttt{entry}. \\ \hline
	fat32-create-dir-entry & entry & Создать массив байт записи в каталоге \texttt{entry}. \\ \hline
	fat32-create-lfn-dir-entries & entry & Создать массив байт записей длинного имени и записи в каталоге \texttt{entry}. \\ \hline
	fat32-find-entry-pos & entry & Найти позицию(смещение в списке блоков . смещение в блоке) для записи в каталоге \texttt{entry} в каталоге \texttt{dir}. \\ \hline
	fat32-get-free-entry-pos & dir-first-block entries-count & Найти позицию (смещение в списке блоков . смещение в блоке) в каталоге, где первый блок \texttt{dir-first-block}, где сможет поместиться количество записей \texttt{entries-count}. При необходимости выделяется новый блок. \\ \hline
	fat32-mark-for-delete & entry & Пометить записей длинного имени и записи в каталоге \texttt{entry} на удаление. \\ \hline
	fat32-add-entry & entry & Добавление записи \texttt{entry} на диск. \\ \hline
	fat32-update-entry & entry changes & Обновить запись в каталоге \texttt{entry} выполнив изменения \texttt{changes}. Если количество записей длинного имени для новой записи отличается, то записи длинного имени и запись в каталоге будет пересоздана. \\ \hline
	is-short-name-sym & symbol & Проверка на символ короткой строки. \\ \hline
	get-trail-num & dir-entries short-name & Получить номер числового хвоста для короткого имени \texttt{short-name}. \\ \hline
	fat32-generate-short-name & dir long-name & Сгенерировать короткое имя из длинного имени \texttt{long-name} в каталоге \texttt{dir}. \\ \hline
	fat32-sort-attributes & attributes is-dir & Отсортировать список атрибутов \texttt{attribute} с проверкой атрибута DIRECTORY в зависимости от того является ли файл каталогом(\texttt{is-dir}). \\ \hline
	fat32-create-special-entries & dir & Создать особые файлы указывающие на текущий каталог и на родительский каталог. \\ \hline
	fat32-is-name-duplicate & dir name & Проверить есть ли в каталоге \texttt{dir} файл с названием \texttt{name}.
\end{xltabular}
\renewcommand{\arraystretch}{1.0} % восстановление сетки

\subsubsection{Класс "<FAT32FileSystem">}

"<FAT32FileSystem"> -- класс файловой системы FAT-32, методы которого являются функциями интерфейса виртуальной файловой системы.

Методы класса "<FAT32FileSystem"> представлены в таблице \ref{FAT32FileSystem:table}

\renewcommand{\arraystretch}{0.8} % уменьшение расстояний до сетки таблицы
\begin{xltabular}{\textwidth}{|p{3.25cm}|p{3.75cm}|>{\setlength{\baselineskip}{0.7\baselineskip}}X|}
	\caption{Описание методов класса "<FAT32FileSystem">\label{FAT32FileSystem:table}}\\
	\hline \centrow \setlength{\baselineskip}{0.7\baselineskip} Название метода & \centrow \setlength{\baselineskip}{0.7\baselineskip} Аргументы & \centrow Описание метода \\
	\hline \centrow 1 & \centrow 2 & \centrow 3\\ \hline
	\endfirsthead
	\caption*{Продолжение таблицы \ref{FAT32FileSystem:table}}\\
	\hline \centrow 1 & \centrow 2 & \centrow 3\\ \hline
	\finishhead
	fs-init & (fat32fs FAT32FileSystem) disk start-sector sector-count &  Загрузка файловой системы FAT32 с диска \texttt{disk}, начиная с сектора \texttt{start-sector}, с количеством секторов \texttt{sector-count}. Чтение блока параметров BIOS, структуры FSInfo, таблицы FAT. \\ \hline
	load-dir & (fat32fs FAT32FileSystem) dir & Раскрыть и сохранить содержимое каталога \texttt{dir}. \\ \hline
	fstat* & (fat32fs FAT32FileSystem) file-hash & Получение метаданных файлового объекта \texttt{file-hash}. \\ \hline
	open-file* & (fat32fs FAT32FileSystem) file-hash & Открыть файл \texttt{file-hash} как поток для чтения и записи. Изменить время последнего доступа. \\ \hline
	new-block & (fat32fs FAT32FileSystem) file-object & Выделить новый блок для файла \texttt{file-object}. \\ \hline
	rename* & (fat32fs FAT32FileSystem) entry new-name & Переименовать файл или каталог \texttt{entry} на \texttt{new-name}. \\ \hline
	remove-file* & (fat32fs FAT32FileSystem) dir entry & Пометить файл \texttt{entry} на удаление. \\ \hline
	remove-dir* & (fat32fs FAT32FileSystem) parent-dir dir & Пометить каталог \texttt{entry} на удаление если он пустой. \\ \hline
	free-space* & (fat32fs FAT32FileSystem) & Получить объём свободного места в текущем разделе в байтах. \\ \hline
	create-file* & (fat32fs FAT32FileSystem) dir name & Создать файл с названием \texttt{name} в каталоге \texttt{dir}. \\ \hline
	create-dir* & (fat32fs FAT32FileSystem) dir name & Создать каталог с названием \texttt{name} в каталоге \texttt{dir}. \\ \hline
	set-attr* & (fat32fs FAT32FileSystem) entry new-attributes & Изменить атрибуты файла \texttt{entry} на \texttt{new-attributes}.
\end{xltabular}
\renewcommand{\arraystretch}{1.0} % восстановление сетки

\subsubsection{Класс "<FAT32File">}

"<FAT32File"> -- класс файла файловой системы FAT-32, методы которого являются функциями интерфейса виртуальной файловой системы.

Методы класса "<FAT32File"> представлены в таблице \ref{FAT32File:table}

\renewcommand{\arraystretch}{0.8} % уменьшение расстояний до сетки таблицы
\begin{xltabular}{\textwidth}{|p{3.25cm}|p{3.75cm}|>{\setlength{\baselineskip}{0.7\baselineskip}}X|}
	\caption{Описание методов класса "<FAT32File">\label{FAT32File:table}}\\
	\hline \centrow \setlength{\baselineskip}{0.7\baselineskip} Название метода & \centrow \setlength{\baselineskip}{0.7\baselineskip} Аргументы & \centrow Описание метода \\
	\hline \centrow 1 & \centrow 2 & \centrow 3\\ \hline
	\endfirsthead
	\caption*{Продолжение таблицы \ref{FAT32File:table}}\\
	\hline \centrow 1 & \centrow 2 & \centrow 3\\ \hline
	\finishhead
	close-file & (self FAT32File) & Закрытие файла. \\ \hline
	tell-file & (self FAT32File) & Получение позиции в файле. \\ \hline
	seek-file & (self FAT32File) offset origin & Сместить позицию в файле на \texttt{offset}, начиная с \texttt{origin}(SET(начало) CUR(текущее) END(конец)). \\ \hline
	read-file & (self FAT32File) size & Прочить \texttt{size} байт из файла и сместить позицию на тоже число. Изменить время последнего доступа. \\ \hline
	write-file & (self FAT32File) buf & Записать в файл массив байт \texttt{buf}, сместив позицию на соответствующее число. Изменить время последнего доступа, время последнего изменения. При необходимости выделить новый блок.
\end{xltabular}
\renewcommand{\arraystretch}{1.0} % восстановление сетки

\subsection{Модули "<iso9660-fs.lsp">, "<iso9660-file.lsp">}

Функции модуля файловой системы ISO-9660(CDFS) представлены в таблице \ref{iso:table}.

\renewcommand{\arraystretch}{0.8} % уменьшение расстояний до сетки таблицы
\begin{xltabular}{\textwidth}{|p{3.5cm}|p{3.5cm}|>{\setlength{\baselineskip}{0.7\baselineskip}}X|}
	\caption{Описание функций, используемых в модуле файловой системы ISO-9660(CDFS)\label{iso:table}}\\
	\hline \centrow \setlength{\baselineskip}{0.7\baselineskip} Название функции & \centrow \setlength{\baselineskip}{0.7\baselineskip} Аргументы & \centrow Описание функции \\
	\hline \centrow 1 & \centrow 2 & \centrow 3\\ \hline
	\endfirsthead
	\caption*{Продолжение таблицы \ref{iso:table}}\\
	\hline \centrow 1 & \centrow 2 & \centrow 3\\ \hline
	\finishhead
	parse-boot-record & Аргументов нет & Парсер оставшейся части структуры дескриптора типа Boot Record. \\ \hline
	parse-primary-desc & Аргументов нет & Парсер оставшейся части структуры дескриптора типа Primary Volume Descriptor. \\ \hline
	parse-set-terminator-desc & Аргументов нет & Парсер оставшейся части структуры дескриптора типа Volume Descriptor Set Terminator. \\ \hline
	parse-descriptor-data & Аргументов нет & Возврат парсера оставшейся части структуры дескриптора типа desc-type. \\ \hline
	parse-date-time-entry & Аргументов нет & Парсер структуры даты и времени в записи в каталоге. \\ \hline
	cdfs-get-date-time-entry & date-time-bytes & Получить список даты и времени из записи в каталоге из массива байт \texttt{date-time-bytes}. \\ \hline
	cdfs-get-attributes-entry & attr-byte & Получить список аттрибутов из записи в каталоге из байта \texttt{attr-byte}. \\ \hline
	parse-cdfs-dir-entry & Аргументов нет & Парсер структуры записи в каталоге. \\ \hline
	parse-descriptor & Аргументов нет & Парсер структура основы дескриптора.
\end{xltabular}
\renewcommand{\arraystretch}{1.0} % восстановление сетки

\subsubsection{Класс "<CDFSFileSystem">}

"<CDFSFileSystem"> -- класс файловой системы ISO-9660(CDFS), методы которого являются функциями интерфейса виртуальной файловой системы.

Методы класса "<CDFSFileSystem"> представлены в таблице \ref{CDFSFileSystem:table}

\renewcommand{\arraystretch}{0.8} % уменьшение расстояний до сетки таблицы
\begin{xltabular}{\textwidth}{|p{3.25cm}|p{3.75cm}|>{\setlength{\baselineskip}{0.7\baselineskip}}X|}
	\caption{Описание методов класса "<CDFSFileSystem">\label{CDFSFileSystem:table}}\\
	\hline \centrow \setlength{\baselineskip}{0.7\baselineskip} Название метода & \centrow \setlength{\baselineskip}{0.7\baselineskip} Аргументы & \centrow Описание метода \\
	\hline \centrow 1 & \centrow 2 & \centrow 3\\ \hline
	\endfirsthead
	\caption*{Продолжение таблицы \ref{CDFSFileSystem:table}}\\
	\hline \centrow 1 & \centrow 2 & \centrow 3\\ \hline
	\finishhead
	fs-init & (self CDFSFileSystem) disk start-sector sector-count & Загрузка файловой системы ISO9660 с диска \texttt{disk}, начиная с сектора \texttt{start-sector}, с количеством секторов \texttt{sector-count}. Чтение основного дескриптора тома, записи в каталоге для корневого каталога из основного дескриптора тома. \\ \hline
	load-dir & (self CDFSFileSystem) dir & Раскрыть и сохранить содержимое каталога \texttt{dir}. \\ \hline
	fstat* & (self CDFSFileSystem) file-hash & Получение метаданных файлового объекта \texttt{file-hash}. \\ \hline
	open-file* & (self CDFSFileSystem) file-hash & Открыть файл как поток для чтения. \\ \hline
	new-block & (self CDFSFileSystem) file-object & Вызов ошибки при попытке выделения нового блока в файловой системе CDFS(ISO9660). \\ \hline
	rename* & (self CDFSFileSystem) entry new-name & Вызов ошибки при попытке переименовании файла или каталога в файловой системе CDFS(ISO9660). \\ \hline
	remove-file* & (self CDFSFileSystem) dir entry & Вызов ошибки при попытке переименовании файла в файловой системе CDFS(ISO9660). \\ \hline
	remove-dir* & (self CDFSFileSystem) parent-dir dir & Вызов ошибки при попытке переименовании каталога в файловой системе CDFS(ISO9660). \\ \hline
	free-space* & (self CDFSFileSystem) & Вызов ошибки при попытке получения свободного места в файловой системе CDFS(ISO9660). \\ \hline
	create-file* & (self CDFSFileSystem) dir name & Вызов ошибки при попытке создании файла в файловой системе CDFS(ISO9660). \\ \hline
	create-dir* & (self CDFSFileSystem) dir name & Вызов ошибки при попытке создании каталога в файловой системе CDFS(ISO9660). \\ \hline
	set-attr* & (self CDFSFileSystem) entry new-attributes & Вызов ошибки при попытке изменения атрибутов файла или каталога в файловой системе CDFS(ISO9660).
\end{xltabular}
\renewcommand{\arraystretch}{1.0} % восстановление сетки

\subsubsection{Класс "<CDFSFile">}

"<CDFSFile"> -- класс файла файловой системы ISO-9660(CDFS), методы которого являются функциями интерфейса виртуальной файловой системы.

Методы класса "<CDFSFile"> представлены в таблице \ref{CDFSFile:table}

\renewcommand{\arraystretch}{0.8} % уменьшение расстояний до сетки таблицы
\begin{xltabular}{\textwidth}{|p{3.25cm}|p{3.75cm}|>{\setlength{\baselineskip}{0.7\baselineskip}}X|}
	\caption{Описание методов класса "<CDFSFile">\label{CDFSFile:table}}\\
	\hline \centrow \setlength{\baselineskip}{0.7\baselineskip} Название метода & \centrow \setlength{\baselineskip}{0.7\baselineskip} Аргументы & \centrow Описание метода \\
	\hline \centrow 1 & \centrow 2 & \centrow 3\\ \hline
	\endfirsthead
	\caption*{Продолжение таблицы \ref{CDFSFile:table}}\\
	\hline \centrow 1 & \centrow 2 & \centrow 3\\ \hline
	\finishhead
	close-file & (self CDFSFile) & Закрытие файла. \\ \hline
	tell-file & (self CDFSFile) & Получение позиции в файле. \\ \hline
	seek-file & (self CDFSFile) offset origin & Сместить позицию в файле на \texttt{offset}, начиная с \texttt{origin}(SET(начало) CUR(текущее) END(конец)). \\ \hline
	read-file & (self CDFSFile) size & Прочитать \texttt{size} байт из файла и сместить позицию на тоже число. \\ \hline
	write-file & (self CDFSFile) buf & Вызов ошибки при попытке записи в файл в файловой системе ISO-9660(CDFS).
\end{xltabular}
\renewcommand{\arraystretch}{1.0} % восстановление сетки

\subsection{Модуль "<fs.lsp">}

\subsubsection{Класс "<FileSystem">}

"<FileSystem"> -- класс абстрактной файловой системы, методы которого являются функциями интерфейса виртуальной файловой системы. Реализации конкретных файловых систем наследуются от этого класса и реализуют его методы.

Методы класса "<FileSystem"> представлены в таблице \ref{FileSystem:table}

\renewcommand{\arraystretch}{0.8} % уменьшение расстояний до сетки таблицы
\begin{xltabular}{\textwidth}{|p{3.25cm}|p{3.75cm}|>{\setlength{\baselineskip}{0.7\baselineskip}}X|}
	\caption{Описание методов класса "<FileSystem">\label{FileSystem:table}}\\
	\hline \centrow \setlength{\baselineskip}{0.7\baselineskip} Название метода & \centrow \setlength{\baselineskip}{0.7\baselineskip} Аргументы & \centrow Описание метода \\
	\hline \centrow 1 & \centrow 2 & \centrow 3\\ \hline
	\endfirsthead
	\caption*{Продолжение таблицы \ref{FileSystem:table}}\\
	\hline \centrow 1 & \centrow 2 & \centrow 3\\ \hline
	\finishhead
	fs-init & (self FileSystem) disk start-sector sector-count & пустой метод, реализация которого находится в конкретной файловой системе. \\ \hline
	load-dir & (self FileSystem) dir & пустой метод, реализация которого находится в конкретной файловой системе. \\ \hline
	fstat* & (self FileSystem) file-hash & пустой метод, реализация которого находится в конкретной файловой системе. \\ \hline
	open-file* & (self FileSystem) file-hash & пустой метод, реализация которого находится в конкретной файловой системе. \\ \hline
	new-block & (self FileSystem) file-object & пустой метод, реализация которого находится в конкретной файловой системе. \\ \hline
	rename* & (self FileSystem) entry new-name & пустой метод, реализация которого находится в конкретной файловой системе. \\ \hline
	remove-file* & (self FileSystem) dir entry & пустой метод, реализация которого находится в конкретной файловой системе. \\ \hline
	remove-dir* & (self FileSystem) parent-dir dir & пустой метод, реализация которого находится в конкретной файловой системе. \\ \hline
	free-space* & (self FileSystem) & пустой метод, реализация которого находится в конкретной файловой системе. \\ \hline
	create-file* & (self FileSystem) dir name & пустой метод, реализация которого находится в конкретной файловой системе. \\ \hline
	create-dir* & (self FileSystem) dir name & пустой метод, реализация которого находится в конкретной файловой системе. \\ \hline
	set-attr* & (self FileSystem) entry new-attributes & пустой метод, реализация которого находится в конкретной файловой системе.
\end{xltabular}
\renewcommand{\arraystretch}{1.0} % восстановление сетки

\subsubsection{Класс "<File">}

"<File"> -- класс абстрактного файла, методы которого являются функциями интерфейса виртуальной файловой системы. Реализации конкретных файлов наследуются от этого класса и в методах реализуют изменение метаданных файла с вызовом метода этого класса.

Методы класса "<File"> представлены в таблице \ref{File:table}

\renewcommand{\arraystretch}{0.8} % уменьшение расстояний до сетки таблицы
\begin{xltabular}{\textwidth}{|p{3.25cm}|p{3.75cm}|>{\setlength{\baselineskip}{0.7\baselineskip}}X|}
	\caption{Описание методов класса "<File">\label{File:table}}\\
	\hline \centrow \setlength{\baselineskip}{0.7\baselineskip} Название метода & \centrow \setlength{\baselineskip}{0.7\baselineskip} Аргументы & \centrow Описание метода \\
	\hline \centrow 1 & \centrow 2 & \centrow 3\\ \hline
	\endfirsthead
	\caption*{Продолжение таблицы \ref{File:table}}\\
	\hline \centrow 1 & \centrow 2 & \centrow 3\\ \hline
	\finishhead
	close-file & (self File) & Закрытие файла. \\ \hline
	tell-file & (self File) & Получение позиции в файле. \\ \hline
	seek-file & (self File) offset origin & Сместить позицию в файле на \texttt{offset}, начиная с \texttt{origin}(SET(начало) CUR(текущее) END(конец)). \\ \hline
	read-file & (self File) size & Прочитать \texttt{size} байт из файла и сместить позицию на тоже число. \\ \hline
	write-file & (self File) buf & Записать в файл массив байт \texttt{buf}, сместив позицию на соответствующее число.
\end{xltabular}
\renewcommand{\arraystretch}{1.0} % восстановление сетки

\subsection{Модуль "<block.lsp">}

Функции модуля работы с блоками диска представлены в таблице \ref{block:table}.

\renewcommand{\arraystretch}{0.8} % уменьшение расстояний до сетки таблицы
\begin{xltabular}{\textwidth}{|p{3.5cm}|p{3.5cm}|>{\setlength{\baselineskip}{0.7\baselineskip}}X|}
	\caption{Описание функций, используемых в модуле работы с блоками\label{block:table}}\\
	\hline \centrow \setlength{\baselineskip}{0.7\baselineskip} Название функции & \centrow \setlength{\baselineskip}{0.7\baselineskip} Аргументы & \centrow Описание функции \\
	\hline \centrow 1 & \centrow 2 & \centrow 3\\ \hline
	\endfirsthead
	\caption*{Продолжение таблицы \ref{block:table}}\\
	\hline \centrow 1 & \centrow 2 & \centrow 3\\ \hline
	\finishhead
	set-block-disk & disk & Установить диск \texttt{disk}, с которого читать блоки. \\ \hline
	set-block-size & size & Установить размер блока диска \texttt{size} в секторах. \\ \hline
	set-block-offset & offset & Установить смещение в секторах для отсчета блоков диска. \\ \hline
	set-blocks-start-num & start-num & Установить значение минимального номера первого блока. \\ \hline
	block-read & num & Прочитать блок с номером \texttt{num}. \\ \hline
	block-write & num buf & Записать буфер buf в блок с номером \texttt{num}. \\ \hline
	get-blocks-pos & blocks offset & Получить индекс блока и смещение по смещению \texttt{offset} в списке блоков \texttt{blocks}. \\ \hline
	get-offset-from-pos & pos & Получить смещение в байтах из позиции(индекс блока . смещение в блоке).
\end{xltabular}
\renewcommand{\arraystretch}{1.0} % восстановление сетки

\subsection{Модуль "<ata.lsp">}

Функции модуля драйвера ATA представлены в таблице \ref{ata:table}.

\renewcommand{\arraystretch}{0.8} % уменьшение расстояний до сетки таблицы
\begin{xltabular}{\textwidth}{|p{3.5cm}|p{3.5cm}|>{\setlength{\baselineskip}{0.7\baselineskip}}X|}
	\caption{Описание функций, используемых в модуле главной загрузочной записи\label{ata:table}}\\
	\hline \centrow \setlength{\baselineskip}{0.7\baselineskip} Название функции & \centrow \setlength{\baselineskip}{0.7\baselineskip} Аргументы & \centrow Описание функции \\
	\hline \centrow 1 & \centrow 2 & \centrow 3\\ \hline
	\endfirsthead
	\caption*{Продолжение таблицы \ref{ata:table}}\\
	\hline \centrow 1 & \centrow 2 & \centrow 3\\ \hline
	\finishhead
	ata-get-status-reg & bit & Получение значения бита \texttt{bit} регистра статуса. \\ \hline
	mk/ata-wait-status & name bit val & Генерация функций ожидания бита регистра статуса. \texttt{name} - имя, \texttt{bit} - номер бита статуса, \texttt{val} - значение для окончания ожидания. \\ \hline
	ata-set-dev & dev & Установка номера устройства. \\ \hline
	ata-set-lba & sec num & Установить номер сектора в режиме LBA. \\ \hline
	ata-check-error & Аргументов нет & Проверка ошибки. \\ \hline
	ata-set-command & cmd & Установить команду контроллера. \\ \hline
	ata-read & size & Чтение данных контроллера, \texttt{size} байт. \\ \hline
	ata-write & arr & Запись массива \texttt{arr} в контроллер. \\ \hline
	ata-identify & Аргументов нет & Чтение служебной информации: число головок, цилиндров, секторов, серийный номер, ... \\ \hline
	ata-read-sectors & dev start num & Читает сектора жесткого диска. \texttt{dev} - устройство: 0 primary master, 1 - slave. \texttt{start} - начальный сектор, \texttt{num} - число секторов. Возвращает массив данных. \\ \hline
	ata-write-sectors & dev start num arr & Пишет сектора жесткого диска. \texttt{dev} - устройство: 0 primary master, 1 - slave. \texttt{start} - начальный сектор, \texttt{num} - число секторов, \texttt{arr} - массив байт данных.
\end{xltabular}
\renewcommand{\arraystretch}{1.0} % восстановление сетки

\subsection{Модульное и интеграционное тестирование разработанной виртуальной файловой системы}

Результаты интеграционного тестирования модуля "<vfs.lsp"> представлены на рисунках \ref{unitIVFS:image} и \ref{unitIVFSdisk:image}

\begin{figure}[H]
	\begin{lstlisting}[language={}]
		"join test"
		"parse path test"
		"load path test"
		"get parent dir test"
		"list dir test"
		"create dir test"
		"remove dir test"
		"change dir test"
		"create file test"
		"remove file test"
		"open file test"
		"is directory test"
		"rename test"
		"free space test"
		"fstat test"
		"work with file test"
		"set attributes test"
		"Успешно завершено тестов: " 26 / 26
		"Все тесты завершены успешно"
	\end{lstlisting}  
	\caption{Результат интеграционного тестирования функций модуля "<vfs.lsp">}
	\label{unitIVFS:image}
\end{figure}

\begin{figure}[H]
	\begin{lstlisting}[language={}]
		"MBR INIT"
		"FAT32 PARTITION INIT"
		"CDFS PARTITION INIT"
		"FAT32 load partition"
		"Create first file named 'first-file.txt'"
		"write to file"
		"read the content of file:" "wrote then read the content of first file"
		"Create dir named 'firstdir'"
		"change-dir to new dir"
		"create second file named 'second file.bin'"
		"write to file"
		"read the content of file:" "wrote then read the content of second created file"
		"CDFS load partition"
		"open first file"
		"content of first file"
		"change-dir to second dir"
		"open second file"
		"first part of second file on first block"
		"second part of second file on second block"
	\end{lstlisting}  
	\caption{Результат интеграционного тестирования модуля "<vfs.lsp">}
	\label{unitIVFSdisk:image}
\end{figure}

Результат интеграционного тестирования модуля "<mbr.lsp"> представлен на рисунке \ref{unitMBR:image}

\begin{figure}[H]
	\begin{lstlisting}[language={}]
		"parse partition rec test"
		"load partition test"
		"Успешно завершено тестов: " 2 / 2
		"Все тесты завершены успешно"
	\end{lstlisting}  
	\caption{Результат интеграционного тестирования функций модуля "<mbr.lsp">}
	\label{unitMBR:image}
\end{figure}

Результаты интеграционного тестирования модулей "<fat32-class.lsp">, "<fat32-fat.lsp">, "<fat32-entry.lsp">, "<fat32-fs.lsp">, "<fat32-file.lsp"> представлены на рисунках \ref{unitFAT:image}, \ref{unitFAT2:image} и \ref{unitFATdisk:image}

\begin{figure}[H]
	\begin{lstlisting}[language={}]
		"parse bios parameter block test"
		"parse fsinfo test"
		"get creation date-time test"
		"get modify date-time test"
		"get access date-time test"
		"get attributes test 1"
		"get attributes test 2"
		"fat elem pos test"
		"get fat chain test"
		"parse lfn test 1"
		"parse lfn test 2"
		"parse dir entry test"
		"fs-init test"
		"load dir test 1"
		"load dir test 2"
		"fstat test"
		"generate creation date-time bytes test"
		"generate modify date-time bytes test"
		"generate access date-time bytes test"
		"generate attribute byte test 1"
		"generate attribute byte test 2"
		"calculate lfn count test 1"
		"calculate lfn count test 2"
		"calculate lfn count test 3"
		"generate checksum test"
		"generate checksum test"
		"create lfn entry test 1"
		"create lfn entry test 2"
		"create lfn entries test"
		"create dir entry test"
		"create both lfn and dir entry test"
		"find entry pos test"
		"update entry test"
		"get free entry pos test"
		"mark entries for delete test"
		"get free block test"
		"update fat elem test1"
		"update fat elem test2"
		"append fat chain test"
		"update entry test1"
		"update entry test2"
		"new block test"
		"generate short name test 1"
		"generate short name test 2"
		"fat32 add entry"
		"rename test"
		"remove file test"
		"remove dir test"
		"free space test"
	\end{lstlisting}  
	\caption{Результат интеграционного тестирования функций модулей "<fat32-class.lsp">, "<fat32-fat.lsp">, "<fat32-entry.lsp">, "<fat32-fs.lsp">, "<fat32-file.lsp">}
	\label{unitFAT:image}
\end{figure}

\begin{figure}[H]
	\begin{lstlisting}[language={}]
		"create file test"
		"create dir test"
		"open file test"
		"read file test"
		"write file test"
		"sort attributes test"
		"set attributes test"
		"Успешно завершено тестов: " 56 / 56
		"Все тесты завершены успешно"
	\end{lstlisting}  
	\caption{Продолжение результатов интеграционного тестирования функций модулей "<fat32-class.lsp">, "<fat32-fat.lsp">, "<fat32-entry.lsp">, "<fat32-fs.lsp">, "<fat32-file.lsp">}
	\label{unitFAT2:image}
\end{figure}

\begin{figure}[H]
	\begin{lstlisting}[language={}]
		"DISK INIT START"
		"DISK INIT END"
		"FS-INIT START"
		"FS-INIT END"
		"EMPTY ROOT DIR LOAD"
		"CREATE FIRST FILE AND FIRST DIR"
		"OPEN FIRST FILE"
		"FIRST FILE WRITE"
		"FIRST FILE READ"
		"first file content"
		"EMPTY FIRST DIR LOAD"
		"CREATE SECOND FILE IN FIRST DIR"
		"OPEN SECOND FILE"
		"SECOND FILE WRITE"
		"SECOND FILE READ"
		"second file content"
		"CREATE SECOND DIR"
		"EMPTY SECOND DIR LOAD"
		"CREATE THIRD FILE IN SECOND DIR"
		"OPEN THIRD FILE"
		"THIRD FILE WRITE"
		"THIRD FILE READ"
		"third file content"
	\end{lstlisting}  
	\caption{Результат интеграционного тестирования модулей "<fat32-class.lsp">, "<fat32-fat.lsp">, "<fat32-entry.lsp">, "<fat32-fs.lsp">, "<fat32-file.lsp">}
	\label{unitFATdisk:image}
\end{figure}

Результаты интеграционного тестирования модулей "<iso9660-fs.lsp">, "<iso9660-file.lsp"> представлены на рисунках \ref{unitISO:image} и \ref{unitISOdisk:image}

\begin{figure}[H]
	\begin{lstlisting}[language={}]
		"get date time from entry test"
		"get attributes from entry test 1"
		"get attributes from entry test 2"
		"get attributes from entry test 3"
		"parse root dir entry test"
		"parse dir entry test 1"
		"parse dir entry test 2"
		"parse dir entry test 3"
		"parse boot record descriptor test"
		"parse descriptor set terminator test"
		"parse primary volume descriptor test"
		"fs init test"
		"load dir test"
		"fstat test"
		"open file test"
		"read file test"
		"Успешно завершено тестов: " 16 / 16
		"Все тесты завершены успешно"
	\end{lstlisting}  
	\caption{Результат интеграционного тестирования функций модулей "<iso9660-fs.lsp">, "<iso9660-file.lsp">}
	\label{unitISO:image}
\end{figure}

\begin{figure}[H]
	\begin{lstlisting}[language={}]
		DISK-MAKE-START
		DISK-MAKE-END
		FS-INIT-START
		FS-INIT-END
		LOAD-ROOT-DIR-START
		LOAD-ROOT-DIR-END
		"content of first file"
		LOAD-DIR-START
		LOAD-DIR-END
		"first part of second file on first block"
		"second part of second file on second block"
	\end{lstlisting}  
	\caption{Результат интеграционного тестирования модулей "<iso9660-fs.lsp">, "<iso9660-file.lsp">}
	\label{unitISOdisk:image}
\end{figure}

Результат интеграционного тестирования модуля "<fs.lsp"> представлен на рисунке \ref{unitVFS:image}

\begin{figure}[H]
	\begin{lstlisting}[language={}]
		"close file test"
		"tell file test"
		"seek file test1"
		"seek file test2"
		"seek file test2.2"
		"seek file test3"
		"read file test1"
		"read file test2"
		"read file test3"
		"write file test1"
		"write file test2"
		"write file test3"
		"write file test4"
		"Успешно завершено тестов: " 13 / 13
		"Все тесты завершены успешно"
	\end{lstlisting}  
	\caption{Результат интеграционного тестирования функций модуля "<fs.lsp">}
	\label{unitVFS:image}
\end{figure}

Результат интеграционного тестирования модуля "<block.lsp"> представлен на рисунке \ref{unitBlock:image}

\begin{figure}[H]
	\begin{lstlisting}[language={}]
		"block set size test"
		"block set offset test"
		"block read test1"
		"block read test2"
		"block write test1"
		"block write test2"
		"get block pos test"
		"get offset from pos test"
		"Успешно завершено тестов: " 8 / 8
		"Все тесты завершены успешно"
	\end{lstlisting}  
	\caption{Результат интеграционного тестирования функций модуля "<block.lsp">}
	\label{unitBlock:image}
\end{figure}

Результат модульного тестирования модуля "<ata.lsp"> представлен на рисунке \ref{unitATA:image}

\begin{figure}[H]
	\begin{lstlisting}[language={}]
		"ata test"
		"Успешно завершено тестов: " 1 / 1
		"Все тесты завершены успешно"
	\end{lstlisting}  
	\caption{Результат интеграционного тестирования функций модуля "<ata.lsp">}
	\label{unitATA:image}
\end{figure}

\subsection{Системное тестирование разработанной виртуальной файловой системы}

Результат системного тестирования разработанной виртуальной файловой системы представлен на рисунке \ref{systemTest:image}

\begin{figure}[H]
	\begin{lstlisting}[language={}]
		> ls
		("boot" "fib.lsp")
		> cd boot
		> cwd
		"/boot"
		> touch test.lsp
		> append test.lsp test-input
		> cat test.lsp
		"test-input"
		> mkdir test-dir
		> copy test.lsp test-dir/test2.lsp
		> cat test-dir/test2.lsp
		"test-input"
		> rm test-dir/test2.lsp
		> rmdir test-dir
		> 
	\end{lstlisting}  
	\caption{Результат системного тестирования разработанной виртуальной файловой системы}
	\label{systemTest:image}
\end{figure}

\begin{comment}
\subsection{Классы, используемые при разработке сайта}

Можно выделить следующий список классов и их методов, использованных при разработке web-приложения (таблица \ref{class:table}). Пример таблицы с уменьшенным межстрочным интервалом.

\renewcommand{\arraystretch}{0.8} % уменьшение расстояний до сетки таблицы
\begin{xltabular}{\textwidth}{|X|p{2.5cm}|>{\setlength{\baselineskip}{0.7\baselineskip}}p{4.85cm}|>{\setlength{\baselineskip}{0.7\baselineskip}}p{4.85cm}|}
\caption{Описание классов Bitrix, используемых в приложении\label{class:table}}\\
\hline \centrow \setlength{\baselineskip}{0.7\baselineskip} Название класса & \centrow \setlength{\baselineskip}{0.7\baselineskip} Модуль, к которому относится класс & \centrow Описание класса & \centrow Методы \\
\hline \centrow 1 & \centrow 2 & \centrow 3 & \centrow 4\\ \hline
\endfirsthead
\caption*{Продолжение таблицы \ref{class:table}}\\
\hline \centrow 1 & \centrow 2 & \centrow 3 & \centrow 4\\ \hline
\finishhead
CMain & Главный модуль & CMain – главный класс страницы web-приложения. После одного из этапов по загрузке страницы в сценарии становится доступным инициализированный системой объект данного класса с именем \$APPLICATION & void ShowTitle(string property\_code = «title», bool strip\_tags = true)
Выводит заголовок страницы
void SetTitle(string title)
Устанавливает заголовок страницы

void ShowCSS(bool external = true, bool XhtmlStyle = true)
Выводит таблицу стилей CSS страницы\\
\hline CFile & Главный модуль & CFile – Класс для работы с файлами и изображениями & array GetFileArray (int file\_id)
Метод возвращает массив, содержащий описание файла (путь к файлу, имя файла, размер) с идентификатором file\_id
\end{xltabular}
\renewcommand{\arraystretch}{1.0} % восстановление сетки

\subsection{Модульное тестирование разработанного web-сайта}

Модульный тест для класса User из модели данных представлен на рисунке \ref{unitUser:image}.

\begin{figure}[ht]
\begin{lstlisting}[language=Python]
from django.test import TestCase
from .models import *
User = get_user_model()


class ShpoTestCases(TestCase):

    def setUp(self) -> None:
        self.user = User.objects.create(username='testtestovich', password='testtestovich', first_name='Sad', last_name='')

    def test_2(self):

        self.assertEqual(self.user.first_name, 'Sad')
        self.assertEqual(self.user.last_name, 'Cat')
        print((self.user))
        print((self.user.first_name))
        print((self.user.last_name))
\end{lstlisting}  
\caption{Модульный тест класса User}
\label{unitUser:image}
\end{figure}

\subsection{Системное тестирование разработанного web-сайта}

На рисунке \ref{main:image} представлена главная страница сайта «Русатом – Аддитивные технологии».
\newpage % при необходимости можно переносить рисунок на новую страницу
\begin{figure}[H] % H - рисунок обязательно здесь, или переносится, оставляя пустоту
\center{\includegraphics[width=1\linewidth]{main1}}
\center{\includegraphics[width=1\linewidth]{main2}}
\center{\includegraphics[width=1\linewidth]{main3}}
\caption{Главная страница сайта «Русатом – Аддитивные технологии»}
\label{main:image}
\end{figure}

На рисунке \ref{menu:image} представлен динамический вывод заголовков, включающий в себя искомые фразы при поиске фраз.

\begin{figure}[ht]
\center{\includegraphics[width=1\linewidth]{menu}}
\caption{Динамический вывод заголовков}
\label{menu:image}
\end{figure}

На рисунке \ref{enter:image} представлен ввод данных для публикации новости.

\begin{figure}[ht]
\center{\includegraphics[width=1\linewidth]{enter}}
\caption{Ввод данных для публикации очень-очень длинной, интересной и полезной новости}
\label{enter:image}
\end{figure}
\end{comment}
